% ============================================================
%  THE ORTYX PROTOCOL
%  Appendix D: Wave-Based Memory Interference Formalism
% ============================================================

\chapter{Wave-Based Memory Interference Formalism}
\label{app:memory}

\noindent
This appendix develops the formal mathematics of the
MEM|8 wave-based memory architecture. The associative
memory efficiency claims that survive the decomposition
at \LevelI{} are derived here.

\section{Memory Wave Functions}
\label{app-d:wave-functions}

Let memory states be represented by complex-valued
wave functions over a representational space $x \in \mathbb{R}^d$
and time $t \ge 0$:
\[
  M_i(x,t) = A_i(x,t)\,e^{i(\omega_i t + \phi_i(x,t))},
\]
where $A_i(x,t) \ge 0$ is a spatially and temporally
varying amplitude envelope, $\omega_i > 0$ is the
characteristic frequency, and $\phi_i(x,t) \in [0, 2\pi)$
is the spatially and temporally varying phase. The global
memory field is:
\[
  \MemField(x,t) = \sum_{i=1}^{N} M_i(x,t).
\]

\section{Retrieval via Resonance}
\label{app-d:retrieval}

Memory retrieval under query function $Q : \mathbb{R}^d \to \mathbb{C}$
is modelled as the resonance score:
\[
  R_i(Q,t) = \left|\int_{\mathbb{R}^d}
  Q(x)\,\overline{M_i(x,t)}\,dx\right|,
\]
where $\overline{M_i}$ denotes complex conjugation. The
memory with highest resonance score is the retrieved
memory.

\begin{proposition}
For orthonormal memory functions $\{M_i\}$ with
$\int M_i \overline{M_j}\,dx = \delta_{ij}$, the retrieval
is exact: $R_i(M_i, t) = \|M_i\|_{L^2}^2$ and
$R_i(M_j, t) = 0$ for $i \neq j$.
\end{proposition}

In practice, memory functions are not orthonormal;
interference between memories affects retrieval fidelity.

\section{Interference Energy and Associative Storage}
\label{app-d:interference}

The interference energy of the global field is:
\[
  E_{\mathrm{int}}(t) = \|\MemField(\cdot, t)\|_{L^2}^2
  = \sum_{i=1}^{N} \|M_i\|_{L^2}^2
  + \sum_{i \neq j} \int M_i(x,t)\,\overline{M_j(x,t)}\,dx.
\]
The cross-term $\sum_{i \neq j} \langle M_i, M_j \rangle$
encodes pairwise associative relationships between memories.

\begin{theorem}
\label{thm:associative-storage}
A superposition of $N$ wave functions encodes $O(N^2)$
pairwise relationships at $O(N)$ storage cost (the $N$
wave functions themselves). Specifically, the cross-term
interference matrix $C_{ij} = \langle M_i, M_j \rangle$
contains $\binom{N}{2}$ distinct pairwise relationships,
all implicitly encoded in the wave superposition without
additional storage.
\end{theorem}

This theorem provides the \LevelI{} grounding for
the associative memory efficiency claim. It follows
from linear algebra and does not depend on the
interpretation of $x$, $\omega_i$, or the biological
correspondence hypothesis.

\section{Amplitude Modulation}
\label{app-d:amplitude}

The emotional amplitude modulation is:
\[
  A_i(x,t) = A_i^0 \exp(-\lambda_i t)\,\bigl(1 + \eta\,e(t)\bigr),
\]
where $\lambda_i > 0$ is the decay constant, $A_i^0$
is the initial amplitude, $e(t) \in \mathbb{R}$ is the
emotional state function, and $\eta > 0$ is the coupling
strength.

\begin{proposition}
Memories with larger coupling $\eta |e(t)|$ decay
more slowly: the effective decay constant under modulation
is approximately $\lambda_i / (1 + \eta |e(t)|)$ for
$\eta |e(t)|$ small. This provides a mathematical
mechanism for emotional enhancement of memory persistence.
\end{proposition}

The operationalization gap noted in Chapters~7 and~12
--- that $e(t)$ lacks a measurement procedure --- remains
open. Proposition above holds for any measurable $e(t)$;
the engineering challenge is specifying $e(t)$ from
observable quantities.

\section{Cross-Modal Binding}
\label{app-d:binding}

Cross-modal binding in the MEM|8 formalism arises from
frequency resonance between memories from different
modalities. Let $M_i^{(a)}$ and $M_j^{(v)}$ denote
auditory and visual memories respectively with
characteristic frequencies $\omega_i^{(a)}$ and
$\omega_j^{(v)}$. Constructive interference occurs when:
\[
  |\omega_i^{(a)} - \omega_j^{(v)}| < \delta_\omega
\]
for a binding threshold $\delta_\omega > 0$. The
binding strength is proportional to the magnitude
of the cross-term $|\langle M_i^{(a)}, M_j^{(v)}\rangle|$,
which is maximized at zero frequency difference.

\begin{proposition}
The binding strength between auditory memory $M_i^{(a)}$
and visual memory $M_j^{(v)}$ decreases monotonically
with $|\omega_i^{(a)} - \omega_j^{(v)}|$ for memories
with otherwise identical structure. The mechanism provides
a mathematical basis for the cross-modal binding claim
at \LevelI{}, contingent on the operationalization of
the characteristic frequencies.
\end{proposition}
